% 随机对象作为可测映射：左为样本空间Omega，右为取值空间E（此处画成实轴）。
% 取值集合B的原像X^{-1}(B)={X in B}可以由不相连的几块组成；可测性要求它属于F，
% 从而P(X in B)有定义。示意图，不含数值。
\begin{tikzpicture}[line width=.6pt,font=\fontsize{8}{10}\selectfont,>=Stealth]
% sample space
\draw[FigureStroke,fill=white,rounded corners=6mm] (0,0) rectangle (42mm,30mm);
\node[FigureInk,anchor=north west] at (2mm,29mm) {$\Omega$};
\fill[FigureBlueFill,draw=FigureBlue] plot[smooth cycle,tension=.8] coordinates {(7mm,6mm) (15mm,4mm) (19mm,10mm) (13mm,15mm) (6mm,12mm)};
\fill[FigureBlueFill,draw=FigureBlue] plot[smooth cycle,tension=.8] coordinates {(26mm,14mm) (34mm,12mm) (37mm,19mm) (30mm,24mm) (24mm,20mm)};
\node[FigureBlue] at (13mm,19mm) {$X^{-1}(B)$};
\draw[FigureBlue,line width=.5pt] (15mm,17.6mm) -- (14mm,13.5mm);
\draw[FigureBlue,line width=.5pt] (18mm,19mm) -- (25mm,19mm);
% points
\coordinate (a) at (12mm,9mm); \coordinate (b) at (31mm,18mm); \coordinate (c) at (22mm,4.5mm);
\foreach \p in {a,b}{\fill[FigureBlue] (\p) circle (.8mm);}
\fill[FigureMuted] (c) circle (.8mm);
\node[FigureMuted,font=\fontsize{7.5}{9}\selectfont,anchor=west] at ($(c)+(1mm,0)$) {$\omega$};
% value space
\begin{scope}[xshift=62mm,yshift=15mm]
 \draw[->,FigureStroke] (0,0) -- (48mm,0) node[anchor=west,FigureInk] {$E=\mathbb R$};
 \fill[FigureYellowFill] (14mm,-2.2mm) rectangle (30mm,2.2mm);
 \draw[FigureYellow,line width=1.1pt] (14mm,0) -- (30mm,0);
 \node[FigureInk] at (22mm,5.5mm) {$B$};
 \coordinate (A) at (18mm,0); \coordinate (Bp) at (26mm,0); \coordinate (C) at (40mm,0);
 \foreach \p in {A,Bp}{\fill[FigureBlue] (\p) circle (.8mm);}
 \fill[FigureMuted] (C) circle (.8mm);
 \node[FigureMuted,font=\fontsize{7.5}{9}\selectfont] at (40mm,4.5mm) {$X(\omega)$};
\end{scope}
\draw[->,FigureBlue,line width=.6pt,shorten >=1.2mm,shorten <=1.2mm] (a) to[out=-10,in=200] (A);
\draw[->,FigureBlue,line width=.6pt,shorten >=1.2mm,shorten <=1.2mm] (b) to[out=10,in=150] (Bp);
\draw[->,FigureMuted,line width=.6pt,dashed,shorten >=1.2mm,shorten <=1.2mm] (c) to[out=-15,in=220] (C);
\node[FigureBlue] at (50mm,24mm) {$X$};
\node[FigureMuted,font=\fontsize{7.5}{9}\selectfont] at (21mm,-5mm) {事件 $\{X\in B\}\in\mathcal F$};
\node[FigureMuted,font=\fontsize{7.5}{9}\selectfont] at (84mm,-5mm) {取值集合 $B\in\mathcal E$};
\end{tikzpicture}
